\section*{Declarações}

\subsection*{Disponibilidade dos dados}

A base experimental é o conjunto de dados publicado por \textcite{DiasLahr2004}, identificado por seu DOI. Os dados derivados deverão distinguir transcrição, cálculos de classificação e hipóteses complementares. A plataforma ReliaBridge é acessível em \parencite{ReliaBridge2026}.

\begin{fillblock}
\noindent\textbf{[A PREENCHER]} Informar o repositório dos dados derivados, scripts e resultados após sua disponibilização. Não declarar acesso aos resultados individuais dos ensaios: o artigo apresenta estatísticas agregadas por espécie.
\end{fillblock}

\subsection*{Conflito de interesse}

Os autores declaram não possuir interesses financeiros concorrentes nem relações pessoais conhecidas que pudessem ter influenciado o trabalho reportado neste artigo.

\subsection*{Financiamento}

\tofill{[Preencher com as agências e números de processo.]}

\subsection*{Declaração de contribuição dos autores (CRediT)}

\tofill{[Preencher conforme a taxonomia CRediT, obrigatória no Engineering Structures: Conceptualization, Methodology, Software, Validation, Formal analysis, Investigation, Resources, Data curation, Writing -- original draft, Writing -- review \& editing, Visualization, Supervision, Project administration, Funding acquisition.]}

\subsection*{Agradecimentos}

\tofill{[Preencher. Incluir agradecimento explícito à fonte dos dados experimentais, caso os autores originais não figurem como coautores.]}
